\documentclass[12pt]{article}
\usepackage{../../ficha}

\begin{document}

\ficha{Potências de Expoente Inteiro e Notação Científica}

\instrucoes{Resolve sem calculadora. Um número está em notação científica quando está na forma $a \times 10^n$, com $1 \le a < 10$ e $n$ inteiro.}

\secao{Potências de Expoente Negativo e Nulo}{Calcula, apresentando o resultado como fração ou número inteiro.}
\begin{exercicios}
  \item $5^0$
  \item $2^{-1}$
  \item $3^{-2}$
  \item $10^{-3}$
  \item $(-2)^{-3}$
  \item $\left(\dfrac{1}{2}\right)^{-3}$
  \item $\left(\dfrac{2}{3}\right)^{-2}$
  \item $\left(-\dfrac{3}{4}\right)^{-1}$
  \item $4^{-2} \times 16$
  \item $2^{-2} + 2^{-1} + 2^0$
\end{exercicios}

\secao{Regras Operatórias}{Escreve na forma de uma única potência.}
\begin{exercicios}
  \item $2^5 \times 2^{-3}$
  \item $3^{-4} \times 3^{-2}$
  \item $5^{2} : 5^{-3}$
  \item $7^{-1} : 7^{4}$
  \item $(2^{-3})^2$
  \item $(10^{-2})^{-3}$
  \item $6^{-3} \times 5^{-3}$
  \item $\left(\dfrac{1}{3}\right)^{4} \times 3^{2}$
  \item $\dfrac{4^{-2} \times 4^{5}}{4^{6}}$
  \item $(-2)^{-4} \times (-2)^{7} : (-2)^{2}$
\end{exercicios}

\secao{Escrever em Notação Científica}{Escreve cada número em notação científica.}
\begin{exercicios}
  \item $45\,000$
  \item $3\,200\,000$
  \item $700\,000\,000$
  \item $0{,}006$
  \item $0{,}000\,42$
  \item $0{,}000\,000\,9$
  \item $230 \times 10^4$
  \item $0{,}05 \times 10^{-2}$
  \item A distância da Terra ao Sol: cerca de $150\,000\,000$\,km.
  \item O diâmetro de um glóbulo vermelho: cerca de $0{,}000\,007$\,m.
\end{exercicios}

\secao{De Notação Científica para Dízima}{Escreve cada número na forma de dízima.}
\begin{exercicios}
  \item $3{,}5 \times 10^3$
  \item $6 \times 10^5$
  \item $1{,}08 \times 10^7$
  \item $4 \times 10^{-2}$
  \item $2{,}7 \times 10^{-4}$
  \item $9{,}01 \times 10^{-6}$
  \item Ordena por ordem crescente: $3 \times 10^{-2}$, $5 \times 10^{-3}$, $2{,}9 \times 10^{-2}$, $1 \times 10^{-1}$.
  \item Qual é maior: $8 \times 10^{5}$ ou $1{,}2 \times 10^{6}$?
\end{exercicios}

\secao{Operações em Notação Científica}{Calcula e apresenta o resultado em notação científica.}
\begin{exercicios}
  \item $(2 \times 10^3) \times (4 \times 10^5)$
  \item $(3 \times 10^{-2}) \times (3 \times 10^{6})$
  \item $(5 \times 10^4) \times (6 \times 10^{2})$
  \item $(8 \times 10^7) : (2 \times 10^3)$
  \item $(1{,}5 \times 10^{-3}) : (5 \times 10^{-6})$
  \item $(2{,}4 \times 10^{5}) : (8 \times 10^{-2})$
  \item $3 \times 10^4 + 5 \times 10^3$
  \item $6{,}2 \times 10^{-3} - 2 \times 10^{-4}$
  \item $(2 \times 10^{-3})^2$
  \item $(3 \times 10^{4})^3$
\end{exercicios}

\secao{Desafios}{Resolve os problemas e apresenta os resultados em notação científica.}
\begin{exercicios}
  \item A luz percorre cerca de $3 \times 10^5$\,km por segundo. Que distância percorre num minuto?
  \item A massa de um átomo de hidrogénio é cerca de $1{,}7 \times 10^{-27}$\,kg. Qual é a massa de $10^{6}$ átomos?
  \item Portugal tem cerca de $1{,}05 \times 10^7$ habitantes e uma área de cerca de $9{,}2 \times 10^4$\,km$^2$. Quantos habitantes há, em média, por km$^2$? (Arredonda às unidades.)
  \item Um vírus mede cerca de $1{,}2 \times 10^{-7}$\,m. Quantos vírus, lado a lado, ocupariam 6\,mm?
  \item Determina $n$ tal que $2^n = \dfrac{1}{32}$.
  \item Determina $n$ tal que $0{,}000\,25 = 2{,}5 \times 10^n$.
  \item Simplifica $\dfrac{(4 \times 10^{-3}) \times (6 \times 10^{8})}{3 \times 10^{2}}$.
\end{exercicios}

\end{document}
